\subsection{Separately concave residuals and minimum cuts}\label{sec:cuts}
The second class removes a large, possibly dense, nonconvex part of the
problem from discretization. Write a rational quadratic as
\begin{equation}\label{eq:cutquad}
F(x)=c_0+\sum_i\bigl(a_ix_i^2+b_ix_i\bigr)+\sum_{i<j}c_{ij}x_ix_j,
\end{equation}
and split the coordinates into a \emph{core} $C$ with $k$ coordinates and a
\emph{residual} $R$. Assume
\begin{equation}\label{eq:cutclass}
a_i\le0\quad(i\in R),\qquad c_{ij}s_is_j\le0\quad(i,j\in R)
\end{equation}
for some signs $s_i\in\{-1,1\}$. Nothing is assumed about core--core or
core--residual coefficients, and the residual interaction graph may be
arbitrary. For a given core, the signs can be found, or shown not to
exist, by propagating relative signs along the residual edges.

\begin{proposition}[Minimum-cut recourse]\label{prop:cut}
Under~\eqref{eq:cutclass}, for every rational core point $v$, every
rational subbox of the residual box (with integer residual coordinates
allowed), and every rational change of the linear coefficients, the
conditional minimum $V(v)=\min_{x_R}F(v,x_R)$ and a minimizing residual
point can be computed by one maximum-flow computation on a graph with
$\abs R+2$ nodes and $O(\abs R+m_R)$ arcs, where $m_R$ is the number of
residual edges. A flow whose value equals the returned cut certifies the
value.
\end{proposition}

\begin{proof}
For fixed $v$ and fixed other residual coordinates, $F$ is concave in each
$x_i$, $i\in R$, so replacing $x_i$ by one of its interval endpoints does
not increase $F$; doing this coordinate by coordinate shows that some
minimizer has all residual coordinates at endpoints. For integer residual
coordinates the endpoints are integers after inward rounding, so the same
holds without enumerating integer labels. Write
$x_i=\alpha_i+\delta_iy_i$ with $y_i\in\{0,1\}$, $\alpha_i$ the lower endpoint
if $s_i=1$ and the upper endpoint if $s_i=-1$, and
$\delta_i=s_i(u_i-\ell_i)$. Substitution gives
$F=K+\sum_ih_iy_i+\sum_{i<j}q_{ij}y_iy_j$ with $q_{ij}=c_{ij}\delta_i\delta_j\le0$,
using $y_i^2=y_i$. With $p_i=h_i+\frac12\sum_{j\ne i}q_{ij}$ and
$r_{ij}=-q_{ij}/2\ge0$, the identity
$q_{ij}y_iy_j=\frac{q_{ij}}2(y_i+y_j)+r_{ij}\abs{y_i-y_j}$ gives
\[
F=K+\sum_ip_iy_i+\sum_{i<j}r_{ij}\abs{y_i-y_j}
=K+\sum_i\min\{0,p_i\}+\cut(y),
\]
where $\cut(y)$ is the capacity of the $s$--$t$ cut that puts $\{i:y_i=1\}$
on the source side in the graph with arcs $i\to t$ of capacity
$\max\{p_i,0\}$, $s\to i$ of capacity $\max\{-p_i,0\}$, and $i\to j$, $j\to i$
of capacity $r_{ij}$. Minimum cuts are computed by maximum flow in strongly
polynomial time, and weak duality certifies optimality
\cite{KolmogorovZabih2004}.
\end{proof}

The class is far from convex. For example, a residual Hessian
$-I-\frac1r\mathbf 1\mathbf 1^\T$ is negative definite with a complete
interaction graph and satisfies~\eqref{eq:cutclass}. The endpoint
reduction and the graph-cut construction are classical; what we use is that
the class is closed under fixing the core, restricting the residual box
and changing linear terms, so that conditional values on corrected core
grids are exact and certified.

By Lemma~\ref{lem:valuefunction}, $V$ has the core curvature bound
$L_C=\max_{i\in C}\max\{2a_i,0\}$ and inherits growth. Treating the core as a
single bag gives the following consequence of Theorem~\ref{thm:valuefn}.

\begin{corollary}[Core search with cut recourse]\label{cor:cutcore}
Assume~\eqref{eq:cutclass} and quadratic growth with constant $g$, and let
$\kappa_C=\max\{1,L_C/g\}$. A certified $2^{-q}$-approximation takes
$f(k,\kappa_C)(I+q+1)^{O(1)}$ bit operations and the exact rational
minimizer takes $f(k,\kappa_C)(I+1)^{O(1)}$. The residual dimension and the
residual interaction graph enter only the polynomial factor.
\end{corollary}

Without growth the core can be searched adaptively with dyadic cells. Take
the core box $[0,1]^k$, cells of side $h_j=2^{-j}$ at level $j$, and the
allowance $e_j=kL_Ch_j^2/8$. Query $V$ at the corners of the current cells,
update the incumbent with the certified completions, keep a cell if the
smallest corner value minus $e_j$ is at most the incumbent value, and split
the kept cells. By Lemma~\ref{lem:interp} applied to $V$ on one cell, the
cell containing a minimizer is always kept, so the smallest kept lower
bound is a certified lower bound, and level
$J=\lceil\frac12\log_2(kL_C/(8\varepsilon))\rceil$ reaches gap $\varepsilon$.
In the worst case the number of kept cells can be exponential in the number
of levels. A random perturbation of the core costs makes it small in
expectation.

\begin{proposition}[Perturbed core costs]\label{prop:smoothed}
Add to $F$ a linear term $\gamma^\T x_C$ whose coefficients are independent
and uniform on $M\ge\max\{2,2^J\}$ equally spaced points of
$[-\sigma,\sigma]$. The search above is correct for every draw, and the
expected number of conditional queries through level $J$ is at most
\[
2^k+8^kJ\Bigl[3+\frac{(1+k/2)L_C}{2\sigma}\Bigr]^k .
\]
\end{proposition}

\begin{proof}
Let $V_\gamma=V+\gamma^\T v$ and $\OPT_\gamma=\min V_\gamma$. A kept cell at
level $j$ has a corner $v$ with $V_\gamma(v)\le U+e_j$, and $U\le\OPT_\gamma+e_j$
because the cell containing a minimizer is queried and
Lemma~\ref{lem:interp} gives a corner within $e_j$ of $\OPT_\gamma$. So every
kept cell has a corner in
$\mathcal N_j=\{v\in h_j\Z^k\cap[0,1]^k:V_\gamma(v)\le\OPT_\gamma+2e_j\}$.
Fix a grid point $v$ and an interior coordinate $i$ ($0<v_i<1$). Comparing
$V_\gamma(v)$ with $V_\gamma(v\pm h_je_i)\ge\OPT_\gamma$ confines $\gamma_i$, if
$v\in\mathcal N_j$, to an interval whose length is
\[
\frac{V(v+h_je_i)+V(v-h_je_i)-2V(v)+4e_j}{h_j}\le L_Ch_j+\frac{4e_j}{h_j}
=L_Ch_j\Bigl(1+\frac k2\Bigr),
\]
by the curvature bound of Lemma~\ref{lem:valuefunction}(a). The interval
depends on $v$ and $V$ but not on the other coordinates of $\gamma$, so these
events are independent across $i$, and each has probability at most
$L_Ch_j(1+k/2)/(2\sigma)+1/M$. Summing the product over the $2^j+1$ values of
each coordinate (two of which are boundary values with probability bound
one) gives
$\E\abs{\mathcal N_j}\le\bigl[2+(2^j-1)\bigl(\frac{L_Ch_j(1+k/2)}{2\sigma}+\frac1M\bigr)\bigr]^k
\le\bigl[3+\frac{(1+k/2)L_C}{2\sigma}\bigr]^k$. A grid point is a corner of at
most $2^k$ cells, each kept cell has $2^k$ children with $2^k$ corners, and
level~$0$ has $2^k$ corners.
\end{proof}

The bound is exponential only in $k$ and in the ratio $L_C/\sigma$ of core
curvature to perturbation size; the residual problem, of any size and
density, enters through the polynomial cost of each query. The optimized
objective is the perturbed one. The same mechanism extends to residual
terms that are polynomial in the core, as long as~\eqref{eq:cutclass} holds
uniformly in the core variables.

\subsection{Why local corrections need exact recourse}\label{sec:local}
The correction of Section~\ref{sec:grids} is global: every coordinate is
rounded, and the allowance $D$ is the sum of all coordinate corrections.
It is tempting to correct only the coordinates of one bag and to use
ordinary grid min-marginals for the rest. That is not valid.

\begin{example}[Star]\label{ex:star}
On $[0,1]^{33}$ let
\[
F(x,y)=x^2-\frac x2\sum_{i=1}^{32}y_i+\sum_{i=1}^{32}y_i^2+\frac{15}{16}x ,
\]
with bags $\{x,y_i\}$ and $L=2$. Minimizing each $y_i$ exactly gives
$y_i=x/4$ and $V(x)=-x^2+\frac{15}{16}x$, so the unique minimizer is
$x=1$, $y_i=1/4$, with value $-1/16$. On the grid $\{0,\frac12,1\}$ the
grid values of $x\mapsto\min_yF$ at $x=0,\frac12,1$ are $0,\frac{23}{32},
\frac{31}{16}$, and the grid incumbent is $0$. The bag cell
$[\frac12,1]\times[0,\frac12]$ of $(x,y_1)$ contains the projection of the
minimizer, yet its smallest corner min-marginal is $\frac{23}{32}$; a bag
correction $\abs BLh^2/8=\frac18$ leaves $\frac{19}{32}>0$, and the cell would
be discarded. The missing error is $32\,\dist(x/4,\{0,\frac12,1\})^2$, which
depends on $x$ and is not removed by normalizing messages.
\end{example}

Exact outside optimization repairs the local rule. For a bag $B$ let
$W_B(v)=\min_{x_{-B}}F(v,x_{-B})$. By Lemma~\ref{lem:valuefunction}, $W_B$
has the same coordinate curvatures, so on a cell of side at most $h$,
\[
\min_{v\text{ corner}}W_B(v)-e_B\le\min_{\text{cell}}W_B,\qquad e_B=\abs BLh^2/8 ,
\]
and only bag coordinates are rounded.

\begin{proposition}[Certified conditional values]\label{prop:interface}
Suppose an oracle returns, at each queried corner $v$, numbers
$\ell_B(v)\le W_B(v)\le u_B(v)$ with $u_B(v)-\ell_B(v)\le\eta$ and a feasible
completion of value $u_B(v)$. Update the incumbent $U$ with all completions
of the current level, and keep a cell if $\min_{\text{corners}}\ell_B-e_B\le
U$. Then every cell containing the projection of a minimizer is kept,
$U\le\OPT+e_B+\eta$, and every kept cell has a corner with
$W_B(v)\le\OPT+2(e_B+\eta)$.
\end{proposition}

\begin{proof}
A cell containing the projection $v^*$ of a minimizer has
$\min_{\text{corners}}\ell_B-e_B\le\min_{\text{corners}}W_B-e_B\le
W_B(v^*)=\OPT\le U$, so it is kept, and its corners give
$U\le\min u_B\le\min W_B+\eta\le\OPT+e_B+\eta$. A kept cell has a corner with
$W_B(v)\le\ell_B(v)+\eta\le U+e_B+\eta$.
\end{proof}

Grid messages do not provide such bounds: in Example~\ref{ex:star} the
grid value overestimates $W_B$ by an amount that varies with the separator
value. Exact conditional values, as in
Sections~\ref{sec:convexrecourse} and~\ref{sec:cuts}, do.
